\section{Applications of the comparison theorems}
\label{section-gaga-applications}

\subsection{Conjugation and Betti numbers}
\label{subsection-gaga-betti-numbers}

Let $\sigma\in\mathop{\rm Aut}(\mathbf C)$. For a scheme $X$ over
$\mathbf C$, write
$$
X^\sigma=X\times_{\mathop{\rm Spec}(\mathbf C),\sigma}
\mathop{\rm Spec}(\mathbf C).
$$

\begin{proposition}
\label{proposition-conjugate-smooth-projective-betti-numbers}
\begin{reference}
\cite[\S 4, no. 18, Proposition 12]{GAGA}
\end{reference}
Let $X$ be a smooth projective algebraic variety over $\mathbf C$. Then $X^{\rm an}$
and $(X^\sigma)^{\rm an}$ have the same Betti numbers.
\end{proposition}

\begin{proof}
For $p,q\geq 0$, flat base change gives a $\sigma$-semilinear
isomorphism
$$
H^q(X,\Omega^p_{X/\mathbf C})
\longrightarrow
H^q(X^\sigma,\Omega^p_{X^\sigma/\mathbf C});
$$
see Cohomology of Schemes, Lemma
\ref{coherent-lemma-flat-base-change-cohomology}. In particular, the two
vector spaces have the same dimension. For a smooth variety, local
coordinates give the canonical comparison
$$
(\Omega^p_{X/\mathbf C})^{\rm an}
\cong \Omega^p_{X^{\rm an}},
$$
where the module on the right is the module of holomorphic $p$-forms. Thus
the cohomology comparison theorem \ref{theorem-gaga-cohomology} identifies
these dimensions with
$$
h^{p,q}(X^{\rm an})
\quad\hbox{and}\quad
h^{p,q}((X^\sigma)^{\rm an}),
$$
respectively, where
$$
h^{p,q}(M)=\dim_{\mathbf C}H^q(M,\Omega^p_M).
$$
Finally, a smooth projective analytification is compact K\"ahler. Hodge
decomposition (using the Dolbeault resolution and the K\"ahler identities),
which is an external analytic input here, gives
$$
b_n(X^{\rm an})=\sum_{p+q=n}h^{p,q}(X^{\rm an})
$$
and the analogous formula for $X^\sigma$. The assertion follows.
\end{proof}

\begin{proposition}[Corollary]
\label{corollary-number-field-embeddings-betti-numbers}
\begin{reference}
\cite[\S 4, no. 18, corollary to Proposition 12]{GAGA}
\end{reference}
Let $V$ be a smooth projective variety over a field $K$ algebraic over
$\mathbf Q$. The Betti
numbers of $V\times_{K,\iota}\mathbf C$ are independent of the embedding
$\iota:K\longrightarrow\mathbf C$.
\end{proposition}

\begin{proof}
The standard extension theorem for field embeddings shows that two
embeddings of $K$ in $\mathbf C$ differ by an automorphism of $\mathbf C$.
Apply Proposition
\ref{proposition-conjugate-smooth-projective-betti-numbers}.
\end{proof}

\begin{remark}[Historical scope]
\label{remark-conjugate-varieties-historical}
Serre observed in \cite[\S 4, no. 18]{GAGA} that a conjugate pair need not be
analytically isomorphic (elliptic curves already show this), and asked there
whether such a pair must be homeomorphic. This historical question is not an
input to Proposition
\ref{proposition-conjugate-smooth-projective-betti-numbers}.
\end{remark}


\subsection{Chow's theorem and algebraicity}
\label{subsection-gaga-chow}

\begin{theorem}[Chow]
\label{theorem-chow-closed-analytic-projective}
\begin{reference}
\cite[\S 4, no. 19, Proposition 13]{GAGA}
\end{reference}
Let $X$ be a projective algebraic variety over $\mathbf C$. Every closed reduced
analytic subspace $Y\subset X^{\rm an}$ is the analytification of a unique
closed reduced subscheme of $X$.
\end{theorem}

\begin{proof}
The vanishing ideal $\mathcal J_Y\subset\mathcal O_{X^{\rm an}}$ is
coherent by Theorem \ref{theorem-oka-cartan-coherence}. By Theorem
\ref{theorem-gaga-essential-surjectivity}, there is a coherent
$\mathcal O_X$-module $\mathcal F$ with
$\mathcal F^{\rm an}\cong\mathcal J_Y$. Compose this isomorphism with the
inclusion into $\mathcal O_{X^{\rm an}}$. The full-faithfulness theorem
\ref{theorem-gaga-fully-faithful} algebraizes the composite to an
$\mathcal O_X$-linear map
$$
u:\mathcal F\longrightarrow\mathcal O_X.
$$
Its image $\mathcal I$ is a coherent ideal. Exactness and faithfulness of
analytification identify $\mathcal I^{\rm an}$ with $\mathcal J_Y$. Hence
the ideal $\mathcal I$ is radical. Indeed, at every closed point
$x\in X(\mathbf C)$, if $a^m\in\mathcal I_x$, then the image of $a$ in
$\mathcal O_{X^{\rm an},x}$ belongs to the radical ideal
$\mathcal J_{Y,x}=\mathcal I_x\mathcal O_{X^{\rm an},x}$. Faithful
flatness in Theorem \ref{theorem-analytification-local-faithfully-flat}
contracts this ideal back to $\mathcal I_x$. Since $X$ is Noetherian and
Jacobson, radicality at all closed stalks implies that $\mathcal I$ is
radical. The closed reduced subscheme defined by $\mathcal I$ therefore has
analytic ideal $\mathcal J_Y$, and hence has analytification $Y$.

If two reduced closed subschemes have analytification $Y$, their ideal
sheaves have equal extensions at every analytic local ring. Faithful
flatness contracts these extensions to equal stalks at every closed point;
coherence and the Jacobson property then identify the ideal sheaves. This
proves uniqueness.
\end{proof}

\begin{proposition}
\label{proposition-compact-analytic-subspace-algebraic}
\begin{reference}
\cite[\S 4, no. 19, Proposition 14]{GAGA}
\end{reference}
Let $X$ be an algebraic variety over $\mathbf C$. Every compact
closed reduced analytic subspace $Y\subset X^{\rm an}$ is algebraic: there
is a unique closed reduced subscheme $Z\subset X$ with
$Z^{\rm an}=Y$.
\end{proposition}

\begin{proof}
Nagata compactification and Chow's lemma (More on Flatness, Theorem
\ref{flat-theorem-nagata} and Limits of Schemes, Lemma
\ref{limits-lemma-chow-finite-type}) give an open immersion
$X\subset\overline X$ with $\overline X$ proper and a proper surjection
$p:X'\to\overline X$ together with an immersion
$X'\to\mathbf P^n_{\mathbf C}$. Replacing $X'$ by its reduction preserves
these properties. The resulting variety $X'$ is proper over $\mathbf C$,
so its immersion is closed and $X'$ is projective. Since $Y$ is compact,
it is closed in the Hausdorff space $\overline X^{\rm an}$. Its analytic structure on
$X^{\rm an}$ glues with the empty closed analytic subspace on
$\overline X^{\rm an}\setminus Y$, and therefore makes $Y$ a closed reduced
analytic subspace of $\overline X^{\rm an}$.

The reduced inverse image
$$
Y'=\left((X')^{\rm an}\times_{\overline X^{\rm an}}Y\right)_{\rm red}
$$
is a closed reduced analytic subspace of $(X')^{\rm an}$. Theorem
\ref{theorem-chow-closed-analytic-projective} algebraizes it to a closed
reduced subscheme $Z'\subset X'$. Let $Z\subset\overline X$ be the reduced
Zariski closure of $p(Z')$. The analytic map sends $(Z')^{\rm an}=Y'$ onto
$Y$. Since $Y$ is closed, Proposition \ref{proposition-image-closure}
shows that the classical point set of $Z^{\rm an}$ is exactly $Y$.
Reduced closed analytic subspaces are determined by their underlying
analytic sets, so $Z^{\rm an}=Y$. All closed points of $Z$ lie in $X$;
the Jacobson property therefore shows that $Z\subset X$.

For uniqueness, if $Z_1,Z_2\subset X$ have analytification $Y$, their ideal
stalks have equal extensions to every $\mathcal O_{X^{\rm an},x}$.
Faithful flatness in Theorem
\ref{theorem-analytification-local-faithfully-flat} contracts them to equal
stalks at every closed point. Coherence and the Jacobson property give
$Z_1=Z_2$.
\end{proof}

\begin{proposition}
\label{proposition-holomorphic-map-from-proper-is-regular}
\begin{reference}
\cite[\S 4, no. 19, Proposition 15]{GAGA}
\end{reference}
Let $X$ be a proper algebraic variety over $\mathbf C$ and let $Y$ be an
algebraic variety over $\mathbf C$. Every holomorphic map
$$
f:X^{\rm an}\longrightarrow Y^{\rm an}
$$
is the analytification of a unique morphism $X\to Y$.
\end{proposition}

\begin{proof}
The graph $\Gamma_f$ is a compact closed analytic subspace of
$(X\times_{\mathbf C}Y)^{\rm an}$. Proposition
\ref{proposition-compact-analytic-subspace-algebraic} makes it the
analytification of a closed reduced subvariety of $X\times_{\mathbf C}Y$.
Proposition \ref{proposition-algebraic-graph-holomorphic-regular} therefore
defines the required algebraic map. For uniqueness, two candidate
maps have closed reduced graphs with the same classical points; the graphs,
and hence the maps, are equal.
\end{proof}

\begin{proposition}[Corollary]
\label{corollary-compact-analytic-space-unique-algebraization}
\begin{reference}
\cite[\S 4, no. 19, corollary to Proposition 15]{GAGA}
\end{reference}
A compact reduced analytic space admits at most one algebraic-variety
structure: any analytic isomorphism between two algebraizations is the
analytification of a unique algebraic isomorphism.
\end{proposition}

\begin{proof}
Every algebraization is proper by Proposition
\ref{proposition-proper-if-and-only-if-compact-analytification}. Apply
Proposition \ref{proposition-holomorphic-map-from-proper-is-regular} to the
analytic isomorphism and to its inverse.
\end{proof}


\subsection{Algebraic and analytic principal bundles}
\label{subsection-gaga-principal-bundles}

\begin{definition}
\label{definition-principal-bundle-analytification-map}
\begin{reference}
\cite[\S 4, no. 20]{GAGA}
\end{reference}
Let $G$ be an algebraic group over $\mathbf C$ and let $X$ be an algebraic
variety over $\mathbf C$. Write
$$
H^1_{\rm Zar}(X,G)
$$
for the pointed set of isomorphism classes of Zariski locally trivial
principal $G$-bundles and
$$
H^1_{\rm an}(X^{\rm an},G^{\rm an})
$$
for the pointed set of isomorphism classes of holomorphically locally
trivial principal $G^{\rm an}$-bundles. Analytification defines a map
$$
\varepsilon_G:H^1_{\rm Zar}(X,G)
\longrightarrow H^1_{\rm an}(X^{\rm an},G^{\rm an}).
$$
This is the torsor interpretation of first cohomology from Cohomology on
Sites, Section \ref{sites-cohomology-section-h1-torsors}.
\end{definition}

\begin{proposition}
\label{proposition-principal-bundle-analytification-injective}
\begin{reference}
\cite[\S 4, no. 20, Proposition 16]{GAGA}
\end{reference}
If $X$ is a proper algebraic variety over $\mathbf C$, then $\varepsilon_G$
is injective. More
precisely, every analytic isomorphism between the analytifications of two
algebraic principal $G$-bundles is algebraic.
\end{proposition}

\begin{proof}
For two algebraic principal $G$-bundles $P$ and $P'$ on $X$, the associated
bundle
$$
\mathop{\rm Isom}\nolimits_G(P,P')
$$
is an algebraic fibre bundle over $X$, and its sections are precisely the
$G$-equivariant isomorphisms $P\to P'$. An analytic isomorphism gives a
holomorphic section of its analytification. Proposition
\ref{proposition-holomorphic-map-from-proper-is-regular} makes the section
algebraic.
\end{proof}

\begin{proposition}
\label{proposition-additive-principal-bundles-gaga}
\begin{reference}
\cite[\S 4, no. 20, Proposition 17]{GAGA}
\end{reference}
If $X$ is a projective algebraic variety over $\mathbf C$ and
$G=\mathbf G_a$, then
$\varepsilon_G$ is bijective.
\end{proposition}

\begin{proof}
Principal $\mathbf G_a$-bundles are torsors under the additive sheaf
$\mathcal O_X$, and hence are classified by $H^1(X,\mathcal O_X)$; see
Cohomology, Lemma \ref{cohomology-lemma-torsors-h1}. The analytic statement
uses $H^1(X^{\rm an},\mathcal O_{X^{\rm an}})$. The comparison isomorphism
of Theorem \ref{theorem-gaga-cohomology} identifies these groups.
\end{proof}

\begin{proposition}
\label{proposition-general-linear-principal-bundles-gaga}
\begin{reference}
\cite[\S 4, no. 20, Proposition 18]{GAGA}
\end{reference}
If $X$ is a projective algebraic variety over $\mathbf C$ and
$G=\mathop{\rm GL}_n$, then $\varepsilon_G$ is bijective.
\end{proposition}

\begin{proof}
Principal $\mathop{\rm GL}_n$-bundles are equivalent to rank $n$ vector
bundles. Let $\mathcal E$ be the locally free coherent analytic module
associated with an analytic vector bundle. By Theorem
\ref{theorem-gaga-essential-surjectivity}, it is isomorphic to
$\mathcal F^{\rm an}$ for a coherent $\mathcal O_X$-module $\mathcal F$.
For every closed point $x\in X(\mathbf C)$, the faithfully flat local map
of Theorem \ref{theorem-analytification-local-faithfully-flat} makes
$$
\mathcal F_x\otimes_{\mathcal O_{X,x}}
\mathcal O_{X^{\rm an},x}
$$
free of rank $n$. Faithfully flat descent for finite projectivity (Algebra,
Proposition \ref{algebra-proposition-ffdescent-finite-projectivity}) shows
that $\mathcal F_x$ is finite projective. Algebra, Theorem
\ref{algebra-theorem-projective-free-over-local-ring} makes it free, and
faithful base change fixes its rank as $n$. Since $X$ is Jacobson and
$\mathcal F$ is coherent, local freeness at every closed point makes
$\mathcal F$ locally free of rank $n$ everywhere. Its frame bundle
algebraizes the given analytic principal bundle. Injectivity is Proposition
\ref{proposition-principal-bundle-analytification-injective}.
\end{proof}

\begin{remark}[Line bundles and the historical formulation]
\label{remark-line-bundles-gaga}
For $n=1$, Proposition
\ref{proposition-general-linear-principal-bundles-gaga} is the isomorphism
$$
\mathop{\rm Pic}(X)\longrightarrow\mathop{\rm Pic}(X^{\rm an}).
$$
Serre's remarks after Proposition 18 in \cite[\S 4, no. 20]{GAGA} express
this in terms of divisor classes on a normal variety and record the earlier
smooth result of Kodaira--Spencer. Those historical attributions are not
additional hypotheses here. Serre also notes that Proposition 18 extends
Kodaira's Theorems 7 and 8 to arbitrary projective varieties, possibly with
singularities, without pursuing those applications.
\end{remark}

\begin{proposition}[Algebraizing a reduction of structure group]
\label{proposition-algebraize-reduction-structure-group}
\begin{reference}
\cite[\S 4, no. 20, Proposition 19]{GAGA}
\end{reference}
Let $H\subset G$ be algebraic groups over $\mathbf C$. Assume that the
quotient $G/H$ exists as an algebraic variety and that
$G\to G/H$ has a rational section. Let $X$ be a proper algebraic variety
over $\mathbf C$ and let $P$ be an analytic principal $H^{\rm an}$-bundle
on $X^{\rm an}$. Then
$P$ is algebraizable if and only if the extended bundle
$$
P\times^H G
$$
is algebraizable.
\end{proposition}

\begin{proof}
A rational section over a nonempty open of $G/H$, and the translates of
its domain by $G$, show that $G\to G/H$ is Zariski locally trivial as an
$H$-bundle.
Necessity is therefore immediate by extension of structure group.

Conversely, choose an algebraic principal $G$-bundle $Q$ and an analytic
isomorphism $Q^{\rm an}\cong P\times^H G$. Form the algebraic associated
bundle
$$
E=Q\times^G(G/H).
$$
The analytic reduction $P$ determines a holomorphic section
$s:X^{\rm an}\to E^{\rm an}$. By Proposition
\ref{proposition-holomorphic-map-from-proper-is-regular}, this section is
algebraic. Pulling the principal $H$-bundle $Q\to E$ back along $s$ gives an
algebraic principal $H$-bundle whose analytification is $P$.
\end{proof}

\begin{definition}
\label{definition-rational-section-condition}
\begin{reference}
\cite[\S 4, no. 20, condition (R)]{GAGA}
\end{reference}
A closed algebraic subgroup $G\subset\mathop{\rm GL}_n$ satisfies
\textup{(R)} if the quotient exists and the projection
$$
\mathop{\rm GL}_n/G\longleftarrow\mathop{\rm GL}_n
$$
has a rational section.
\end{definition}

\begin{proposition}
\label{proposition-rational-section-principal-bundles-gaga}
\begin{reference}
\cite[\S 4, no. 20, Proposition 20]{GAGA}
\end{reference}
Let $G\subset\mathop{\rm GL}_n$ satisfy \textup{(R)}. For every projective
algebraic variety $X$ over $\mathbf C$, the map
$$
\varepsilon_G:H^1_{\rm Zar}(X,G)
\longrightarrow H^1_{\rm an}(X^{\rm an},G^{\rm an})
$$
is bijective.
\end{proposition}

\begin{proof}
Given an analytic principal $G$-bundle $P$, Proposition
\ref{proposition-general-linear-principal-bundles-gaga} algebraizes
$P\times^G\mathop{\rm GL}_n$. Proposition
\ref{proposition-algebraize-reduction-structure-group}, with
$H=G$ and the ambient group $\mathop{\rm GL}_n$, then algebraizes $P$.
Injectivity is Proposition
\ref{proposition-principal-bundle-analytification-injective}.
\end{proof}

\begin{remark}[Examples and historical boundary]
\label{remark-rational-section-condition-examples}
Condition \textup{(R)} holds for solvable subgroups by Rosenlicht's theorem,
for $\mathop{\rm SL}_n$ via the determinant section, and for
$\mathop{\rm Sp}_{2m}$ by rational symplectic Gram--Schmidt. Serre also
records in \cite[\S 4, no. 20]{GAGA} the then-open simply connected
semisimple conjecture and the failure of \textup{(R)} for the unimodular
orthogonal group in dimensions at least $3$; in the latter case he leaves
bijectivity of $\varepsilon_G$ open. These historical observations are not
used in Proposition
\ref{proposition-rational-section-principal-bundles-gaga}.
\end{remark}
